% ===================== PREAMBULE COMMUN (à copier VERBATIM en tête de CHAQUE .tex) =====================
\documentclass[11pt,a4paper]{article}
\usepackage[utf8]{inputenc}
\usepackage[T1]{fontenc}
\usepackage{lmodern}
\usepackage[french]{babel}
\usepackage{amsmath,amssymb,mathtools}
\usepackage{icomma}
\usepackage{xcolor}
\usepackage{colortbl}
\usepackage{tikz}
\usepackage{pgfplots}
\usepackage[most]{tcolorbox}
\usepackage{enumitem}
\usepackage{array}
\usepackage{booktabs}
\usepackage{multicol}
\usepackage{fancyhdr}
\usepackage{caption}
\captionsetup{labelformat=empty,justification=centering,font=small}
\usepackage[hidelinks]{hyperref}
\usetikzlibrary{arrows.meta,positioning,calc,angles,quotes,patterns,backgrounds,shapes.geometric}
\pgfplotsset{compat=1.18}
\usepackage[a4paper,margin=2.2cm,headheight=26pt,headsep=10pt]{geometry}

\definecolor{primary}{RGB}{219,54,77}
\definecolor{primarydark}{RGB}{150,28,46}
\definecolor{primarylight}{RGB}{250,224,228}
\definecolor{defcol}{RGB}{0,114,178}
\definecolor{thmcol}{RGB}{0,140,120}
\definecolor{excol}{RGB}{0,135,90}
\definecolor{methcol}{RGB}{90,90,100}
\definecolor{courbe}{RGB}{40,90,200}
\definecolor{courbeB}{RGB}{210,60,40}
\definecolor{courbeC}{RGB}{30,150,80}
\definecolor{axiscol}{RGB}{60,60,60}
\definecolor{gridcol}{RGB}{200,205,215}

\tcbset{boxbase/.style={enhanced,breakable,boxrule=0.9pt,arc=2.5pt,
   left=3mm,right=3mm,top=2.3mm,bottom=2.3mm,
   fonttitle=\bfseries,coltitle=white,toptitle=0.6mm,bottomtitle=0.6mm}}
\newtcolorbox{methode}[1][]{boxbase,colback=methcol!7,colframe=methcol,sharp corners,
  title={\textsc{Méthode}\ifx\relax#1\relax\relax\else\textmd{ --- \itshape #1}\fi}}
\newtcolorbox{exemple}[1][]{boxbase,colback=excol!5,colframe=excol,
  title={\textsc{Exemple}\ifx\relax#1\relax\relax\else\textmd{ : \itshape #1}\fi}}
\newtcolorbox{idee}[1][]{boxbase,colback=defcol!6,colframe=defcol,
  title={\textsc{L'essentiel}\ifx\relax#1\relax\relax\else\textmd{ : \itshape #1}\fi}}
\newtcolorbox{piege}[1][]{boxbase,colback=primary!7,colframe=primary,
  title={\textsc{Piège}\ifx\relax#1\relax\relax\else\textmd{ : \itshape #1}\fi}}

\pgfplotsset{repere/.style={axis lines=middle,
    axis line style={-{Stealth[length=2.2mm]},axiscol,line width=0.8pt},
    label style={font=\footnotesize},tick label style={font=\scriptsize},
    every axis x label/.style={at={(current axis.right of origin)},anchor=north west},
    every axis y label/.style={at={(current axis.above origin)},anchor=south east},
    xlabel={$x$},ylabel={$y$},grid=both,minor tick num=0,
    major grid style={gridcol,line width=0.4pt},samples=200,clip=false}}

\pagestyle{fancy}\fancyhf{}
\fancyhead[L]{\small\textcolor{primarydark}{\textbf{Fiche 4} — Les fonctions}}
\fancyhead[R]{\small\textcolor{primarydark}{\textsc{Carnet d'entraînement}}}
\fancyfoot[C]{\small\thepage}
\renewcommand{\headrulewidth}{0.8pt}
\renewcommand{\headrule}{\hbox to\headwidth{\color{primary}\leaders\hrule height \headrulewidth\hfill}}
\usepackage{titlesec}
\titleformat{\section}{\Large\bfseries\color{primarydark}}{\thesection}{0.6em}{}
\titleformat{\subsection}{\large\bfseries\color{primary}}{\thesubsection}{0.6em}{}

\newcommand{\R}{\mathbb{R}}\newcommand{\Z}{\mathbb{Z}}\newcommand{\N}{\mathbb{N}}
\newcommand{\Df}{D_f}
\newcommand{\croit}{\textcolor{thmcol}{\Large$\nearrow$}}
\newcommand{\decroit}{\textcolor{thmcol}{\Large$\searrow$}}

\newcommand{\exercice}[1]{\medskip\noindent{\color{primarydark}\bfseries\large Exercice #1}\par\nopagebreak\smallskip}
\newcommand{\titrecarnet}[3]{%
\begin{center}\begin{tikzpicture}
\node[rectangle,rounded corners=7pt,fill=primary,text=white,minimum width=15cm,
  minimum height=1.7cm,align=center,inner sep=8pt]
  {{\LARGE\bfseries #1}\\[2pt]{\normalsize #2}};
\end{tikzpicture}\end{center}
\begin{center}\itshape\color{primarydark}#3\end{center}\medskip}

% ---- RÈGLES D'USAGE DES MACROS ----
% * \croit et \decroit s'emploient en mode TEXTE (dans une cellule de tableau),
%   JAMAIS entre $...$ (ils contiennent déjà un $\nearrow$).
% * \titrecarnet{Titre}{Sous-titre}{Ligne en italique} : bandeau de titre.
% * \exercice{1} : titre d'exercice.
% * Boîtes : \begin{idee}...\end{idee}, \begin{exemple}[titre]...\end{exemple},
%   \begin{methode}[titre]...\end{methode}, \begin{piege}[titre]...\end{piege}.
% Le corps du document commence après \begin{document}.
\begin{document}
\titrecarnet{Thème 2 — Parité, symétries \& périodicité}{Corrigé — niveau Difficile}{Détail des calculs.}

\exercice{1}
\begin{enumerate}[label=\alph*),leftmargin=1.8em,topsep=2pt]
  \item $f(-x)=(-x)^4-2(-x)^2+1=x^4-2x^2+1=f(x)$ : $f$ est \textbf{paire}.
  \item Une fonction paire est symétrique par rapport à l'axe des ordonnées : l'\textbf{axe de symétrie est $x=0$}.
  \item $f(x)=(x^2-1)^2=0\Leftrightarrow x^2-1=0\Leftrightarrow x=\pm 1$. Racines : $\boxed{x=-1 \text{ et } x=1}$ (chacune double).
\end{enumerate}

\exercice{2}
$f(x)=x^2+4x+7$, $a=1$, $b=4$.
\begin{enumerate}[label=\alph*),leftmargin=1.8em,topsep=2pt]
  \item Axe : $x=-\dfrac{b}{2a}=-\dfrac{4}{2}=\boxed{-2}$.
  \item Forme canonique : $x^2+4x+7=(x+2)^2-4+7=(x+2)^2+3$. Donc $f(x)=(x-(-2))^2+3$, sommet $\boxed{(-2;3)}$.
        (Cohérent : abscisse du sommet $=-2=$ axe de symétrie.)
\end{enumerate}

\exercice{3}
\begin{enumerate}[label=\alph*),leftmargin=1.8em,topsep=2pt]
  \item On factorise la différence de deux carrés :
        \[ \cos^4 x-\sin^4 x=(\cos^2 x-\sin^2 x)(\cos^2 x+\sin^2 x). \]
        Or $\cos^2 x+\sin^2 x=1$ et $\cos^2 x-\sin^2 x=\cos(2x)$ (formule de duplication). Donc $g(x)=\cos(2x)$. \checkmark
  \item $\cos(2x)$ a pour période $\dfrac{2\pi}{2}=\pi$. Donc $\boxed{T=\pi}$.
\end{enumerate}

\exercice{4}
\begin{enumerate}[label=\alph*),leftmargin=1.8em,topsep=2pt]
  \item Formule d'addition : $\sqrt2\,\sin\!\Big(x+\dfrac{\pi}{4}\Big)=\sqrt2\Big(\sin x\cos\dfrac{\pi}{4}+\cos x\sin\dfrac{\pi}{4}\Big)=\sqrt2\Big(\sin x\cdot\dfrac{\sqrt2}{2}+\cos x\cdot\dfrac{\sqrt2}{2}\Big)=\sin x+\cos x=f(x)$. \checkmark
  \item $\sin\!\Big(x+\dfrac{\pi}{4}\Big)$ a pour période $2\pi$ ; le facteur $\sqrt2$ et le décalage $\dfrac{\pi}{4}$ ne changent pas la période. Donc $\boxed{T=2\pi}$.
  \item On teste en $x=\dfrac{\pi}{4}$ : $f\!\Big(\dfrac{\pi}{4}\Big)=\sqrt2\,\sin\!\Big(\dfrac{\pi}{2}\Big)=\sqrt2$, tandis que $f\!\Big(-\dfrac{\pi}{4}\Big)=\sqrt2\,\sin(0)=0$.
        Comme $f(-\tfrac{\pi}{4})=0$ vaut ni $f(\tfrac{\pi}{4})=\sqrt2$ ni $-f(\tfrac{\pi}{4})=-\sqrt2$, la fonction $f$ n'est \textbf{ni paire ni impaire}.
\end{enumerate}

\exercice{5}
\begin{enumerate}[label=\alph*),leftmargin=1.8em,topsep=2pt]
  \item $f(-x)=\lvert\sin(-x)\rvert=\lvert-\sin x\rvert=\lvert\sin x\rvert=f(x)$ : $f$ est \textbf{paire}.
  \item $f(x+\pi)=\lvert\sin(x+\pi)\rvert=\lvert-\sin x\rvert=\lvert\sin x\rvert=f(x)$ : $\pi$ est une période. Elle est plus petite que $2\pi$ : la valeur absolue « replie » les arches négatives du sinus vers le haut, ce qui double la fréquence apparente. Donc $\boxed{T=\pi}$.
\end{enumerate}

\begin{center}
\begin{tikzpicture}
\begin{axis}[repere,width=13cm,height=4cm,xmin=-1.7,xmax=7,ymin=-0.4,ymax=1.4,
   xtick={3.1416,6.283},xticklabels={$\pi$,$2\pi$},ytick={1},samples=200,
   title={\footnotesize $f(x)=|\sin x|$ — période $\pi$}]
  \addplot[courbe,line width=1.2pt,domain=-1.6:7]{abs(sin(deg(x)))};
\end{axis}
\end{tikzpicture}
\end{center}
\end{document}
